\section{Experiments}

\subsection{Dataset}

We validate the framework through systematic evaluation of 100 short stories across three quality categories designed to test discriminative capacity. The stratified corpus enables empirical testing of whether LLM-based evaluation can distinguish meaningful quality differences while controlling for potential confounds including genre, period, and provenance. To control for length effects, all stories fall within the 2,000--8,000 word range characteristic of the short fiction form.

\subsubsection{Canonical Literature Corpus}

The canonical subset comprises 50 short stories representing works consistently recognized by academic institutions, literary anthologies, and critical scholarship. Authors include Chekhov, Kafka, Joyce, Faulkner, Borges, Lu Xun, García Márquez, Hemingway, Poe, Woolf, and others whose works appear regularly in university syllabi and literary criticism. Publication dates span 1835 to 1990, encompassing major literary movements including Romanticism, Realism, Modernism, and Postmodernism. Geographic diversity includes Russian, European, North American, Latin American, and Japanese literary traditions. This heterogeneity ensures the framework encounters varied narrative techniques, thematic concerns, and cultural contexts characteristic of the literary canon.

\subsubsection{Pulp Fiction Corpus}

The pulp fiction subset comprises 30 short stories published between 1880 and 1950 in commercial genre magazines. Works drawn from publications including \textit{All-Story}, \textit{Weird Tales}, and \textit{Black Mask} represent adventure pulp, horror pulp, western pulp, detective fiction, and early science fiction. Authors include Burroughs, Lovecraft, Grey, and others writing for mass-market periodicals. This corpus provides a controlled comparison category distinguished from canonical literature by critical reception, institutional positioning, and commercial publication context. Pulp fiction exhibits technical narrative competence and reader engagement while typically employing genre conventions, formulaic plot structures, and commercial narrative strategies.

\subsubsection{LLM-Generated Story Corpus}

The LLM-generated subset comprises 20 short stories selected from the lars76/story-evaluation-llm dataset on Hugging Face~\cite{huggingface-story-eval}, which contains narratives generated by various open-source LLM models including Mistral variants, Gemma variants, and Llama variants. The corpus exhibits intentional quality stratification based on human evaluation scores from the source dataset: 10 stories received high-quality ratings (overall scores 3.58--3.78 on 5-point scale), 5 stories exhibit average quality (overall score 3.40), and 5 stories demonstrate low quality (overall scores 2.56--2.58). This stratified distribution enables testing whether LLM-based evaluation can discriminate quality levels within AI-generated content and identify distinctive characteristics of synthetic narratives across different generative models.

\begin{table}[t]
\centering
\caption{Evaluation corpus composition and distribution}
\label{tab:corpus_composition}
\small
\begin{tabular}{lccl}
\toprule
\textbf{Category} & \textbf{Stories} & \textbf{Evaluations} & \textbf{Quality Characteristics} \\
\midrule
Canonical & 50 & 300 & Academically recognized, \\
          &    &     & regularly anthologized \\
\midrule
Pulp Fiction & 30 & 180 & Commercial periodicals, \\
             &    &     & genre conventions \\
\midrule
LLM-Generated & 20 & 120 & Multi-model generated, \\
              &    &     & stratified quality \\
\midrule
\textbf{Total} & \textbf{100} & \textbf{600} & \textbf{3 layers $\times$ 2 modes} \\
\bottomrule
\end{tabular}
\vspace{0.1cm}

\small
\textit{Notes:} Each story evaluated across 3 analytical layers (Cultural, Emotional, Existential) in 2 modes (content-limit, title-limit). Total: 100 stories $\times$ 3 layers $\times$ 2 modes = 600 evaluations. LLM-generated stories include high (n=10), average (n=5), and low (n=5) quality tiers.
\end{table}

The three-category design provides complementary quality contrasts. Comparing canonical literature to pulp fiction tests whether the framework discriminates works distinguished by critical reception and institutional canonization. Comparing human-authored stories (canonical and pulp) to LLM-generated narratives tests whether the framework identifies distinctive characteristics of AI-generated content. Comparing quality tiers within LLM-generated stories tests whether the framework tracks deliberate quality variation. Together these comparisons provide converging evidence for discriminative validity.

\subsection{Evaluation Procedure}

The experimental validation encompasses three analytical layers: Layer 4 (Cultural Representation), Layer 5 (Emotional-Psychological Representation), and Layer 6 (Existential-Philosophical Representation). As described in Methodology, each story undergoes evaluation in two modes (content-limit and title-limit) by two evaluators (iterative evaluator with 5-round reflection and independent validator). Table~\ref{tab:evaluation_matrix} summarizes the complete evaluation matrix.

\begin{table}[t]
\centering
\caption{Complete evaluation matrix across three SAGE layers}
\label{tab:evaluation_matrix}
\small
\begin{tabular}{lcccc}
\toprule
\textbf{Layer} & \textbf{Dimensions} & \textbf{Stories} & \textbf{Modes} & \textbf{Evaluations} \\
\midrule
Layer 4 (Cultural) & 4 & 100 & 2 & 200 \\
Layer 5 (Emotional) & 4 & 100 & 2 & 200 \\
Layer 6 (Existential) & 4 & 100 & 2 & 200 \\
\midrule
\textbf{Total} & \textbf{12} & \textbf{100} & \textbf{2} & \textbf{600} \\
\bottomrule
\end{tabular}
\vspace{0.1cm}

\small
\textit{Notes:} Each evaluation includes iterative assessment (5 rounds) + independent validation (1 round). All three layers evaluate identical story set. Success rate: 600/600 (100\%). Dimensions: L4 (Intersectional Power Dynamics, Cultural Voice \& Perspective, Cultural Specificity, Cultural Pattern Complexity); L5 (Affective Complexity, Psychological Interiority, Emotional Granularity, Emotional-Narrative Coherence); L6 (Life Philosophy, Moral Reflection, Human Condition, Meaning Exploration).
\end{table}

\subsection{Implementation Details}

All LLM-based evaluations employ GPT-5-mini (OpenAI) with reasoning effort parameter set to ``medium'' to balance analytical depth and response consistency. The model supports maximum output tokens of 128,000, sufficient for processing complete short stories with detailed evaluation prompts. Each evaluation produces structured JSON outputs parsed to extract dimension scores, confidence ratings, textual evidence citations, and reasoning summaries.

\subsubsection{Prompt Engineering}

Prompts incorporate explicit theoretical frameworks as detailed in Methodology. For Layer 4 (Cultural Representation), prompts reference Bourdieu's field theory, Said's Orientalism, Geertz's thick description, and structural anthropology. For Layer 5 (Emotional-Psychological Representation), prompts reference affect theory (Sedgwick, Berlant, Ngai), narrative psychology (Cohn, Bruner, Bakhtin), emotion granularity research (Barrett), and New Critical concepts of organic unity. For Layer 6 (Existential-Philosophical Representation), prompts reference existentialist philosophy, moral philosophy, philosophical anthropology, and hermeneutic frameworks as specified in the theoretical foundation.

Prompts establish core evaluation principles including evidence sovereignty (all scores must cite textual evidence), strict layer boundaries (evaluate only target layer dimensions), and theoretical consistency (apply frameworks uniformly across texts). Prompts explicitly instruct models to avoid common pitfalls: conflating explicitness with depth, imposing evaluator's emotional or cultural norms on texts, crossing layer boundaries, and making inferences unsupported by evidence.

\subsubsection{Example Prompts}

\begin{promptbox}[title={\small Layer 4, Content-Limit Mode, Round 1}]
\small
You are evaluating the cultural representation quality of a literary work based solely on textual evidence. You have been provided with the story text but no information about the title, author, or publication context.

Your task is to assess four dimensions of cultural representation:

\textbf{Intersectional Power Dynamics (IPD):} Grounded in Bourdieu's theory of cultural capital alongside Marxist, feminist, and postcolonial frameworks, evaluate complexity and configuration of power relations as they operate through social structures. Power may be organized along various social axes including class, gender, race/ethnicity, kinship, occupation, age, religion, institutional position. Focus on: how power is distributed, negotiated, and legitimated through social structures; interaction between forms of capital (economic, cultural, symbolic); structural mechanisms (institutions, norms, hierarchies); complexity from simple binary oppositions to multi-layered structural tensions.

[Additional dimensions: CVP, CSP, CPC with similar theoretical grounding...]

\textbf{CRITICAL:} Base your evaluation ONLY on what is present in the text itself. Do not use any prior knowledge about literary works, authors, or cultural contexts beyond what the text explicitly provides. Cite specific passages to support all claims.

Provide scores (1.0--5.0) and confidence ratings (1--5) for each dimension with detailed reasoning.
\end{promptbox}

% \vspace{0.2cm}

\begin{promptbox}[title={\small Layer 4, Title-Limit Mode, Round 1}]
\small
You are conducting a meta-critical analysis of a literary work's cultural representation quality based on scholarly consensus and critical reception.

You have been provided with the work's title and author but NOT the text itself. Your task is to synthesize assessments from literary critics, scholars, and academic analyses of this specific work.

For each of the four cultural dimensions (IPD, CVP, CSP, CPC), draw on your knowledge of how critics and scholars have evaluated this work. Consider:

\textbf{Critical Consensus:} What do major literary critics and scholars say about this work's treatment of power dynamics, cultural positioning, cultural specificity, and structural sophistication? Cite specific critics or critical schools where relevant.

\textbf{Literary Historical Context:} How has this work's cultural representation been assessed across different critical periods? Has critical evaluation evolved over time?

\textbf{Scholarly Reliability:} Distinguish between well-grounded critical consensus supported by extensive analysis and more speculative or partial interpretations.

\textbf{CRITICAL:} Base your assessment on the weight of scholarly consensus. If critical opinion is divided, note the division and explain the grounds of disagreement. Do not impose personal judgments where scholarly opinion provides clear guidance.

Provide dimension scores (1.0--5.0) based on critical consensus with detailed reasoning explaining which scholarly perspectives informed your assessment.
\end{promptbox}

\begin{table}[t]
\centering
\caption{Implementation specifications}
\label{tab:implementation_specs}
\small
\resizebox{\columnwidth}{!}{%
\begin{tabular}{lc}
\hline
\textbf{Parameter} & \textbf{Value} \\
\hline
\multicolumn{2}{l}{\textit{LLM Configuration}} \\
Model & GPT-5-mini \\
Provider & OpenAI \\
Reasoning effort & medium \\
Max output tokens & 128,000 \\
Output format & Structured JSON \\
\hline
\multicolumn{2}{l}{\textit{Scoring Configuration}} \\
Dimension score range & 1.0--5.0 \\
Confidence rating range & 1--5 \\
Convergence threshold & $< 0.3$ (Round 4$\rightarrow$5) \\
Inter-rater agreement threshold & MAD $< 0.5$ \\
\hline
\multicolumn{2}{l}{\textit{Evaluation Structure}} \\
Iterative evaluator rounds & 5 \\
Independent validator rounds & 1 \\
\hline
\end{tabular}%
}
\end{table}

\subsection{Statistical Analysis}

We analyze evaluation reliability, discriminative validity, and dimensional independence through multiple complementary analyses organized around the four research questions.

For evaluation reliability (RQ1), we track score trajectories across the five iterative rounds for each dimension and story. Dimensions exhibiting minimal score change (absolute difference $< 0.3$) from Round 4 to Round 5 are classified as convergent. We calculate convergence rates by dimension, layer, mode, and story category. Inter-rater reliability between iterative evaluator and independent validator is quantified through mean absolute difference (MAD), with agreement rates meeting the MAD $< 0.5$ threshold indicating acceptable reliability.

For discriminative validity (RQ2), we employ independent samples t-tests to compare mean scores across the three story categories (canonical vs.\ pulp vs.\ LLM-generated) for each layer. Effect sizes are reported through Cohen's $d$ with standard interpretation thresholds: small ($d$=0.2), medium ($d$=0.5), large ($d$=0.8). One-way ANOVA tests assess overall genre effects, followed by post-hoc pairwise comparisons with Bonferroni correction. Dimension-level analysis identifies which specific dimensions contribute most to category discrimination.

For evaluation robustness (RQ3), we compare scores between content-limit and title-limit modes using paired samples t-tests. Mode effects reveal whether LLM evaluations rely primarily on textual analysis, memorized critical consensus, or both sources.

For dimensional independence (RQ4), we compute Pearson and Spearman correlation coefficients between all pairs of Layer 4, Layer 5, and Layer 6 scores. Moderate correlations ($r < 0.7$) support dimensional independence, while high correlations would suggest substantial overlap requiring framework revision. We also compare cross-layer score distributions to test whether the three analytical dimensions exhibit differential genre sensitivity.

The evaluation framework produces both quantitative scores enabling statistical analysis and qualitative rationales permitting examination of how LLM assessment responds to specific literary features. This dual output supports both hypothesis testing and interpretive analysis of evaluation patterns.